\newpage\handout
{Drill Problems: Week 2.3}
{\textsc{Scholastic Aptitude Test (SAT)}}
{\href{https://creativecommons.org/licenses/by-nc-sa/4.0/}{CC BY-NC-SA 4.0 license}}
{Author: \BookAuthor}{Release: \generatedOn}
{SAT: Drill Problems (2.3)}
% Copyright & Disclaimer
\begin{center}
  \begin{minipage}{0.85\textwidth}
    {\small\textbf{Purpose and Usage:}}\\[0.2cm]
    {\footnotesize
    This material has been developed for internal training and educational 
    purposes at Hans edu LLC. It is intended for use within our organization 
    and should not be distributed, sold, or used for commercial purposes 
    outside of our educational programs.}\\[0.5cm]
    
    {\small\textbf{For Our Community:}}\\[0.2cm]
    {\footnotesize
    Students and staff are welcome to use this material in their studies and 
    teaching at Hans edu LLC. While we encourage active engagement with the 
    content, please respect that this is proprietary material. Any 
    reproduction or distribution outside of our organization's educational 
    activities is not permitted.}\\[0.5cm]
    
    {\small\textbf{Content and Attribution:}}\\[0.2cm]
    {\footnotesize
    This material represents our adaptation of various established mathematics 
    textbooks, reorganized and enhanced for our teaching context. While we've 
    added our own pedagogical improvements, we maintain proper attribution to 
    original sources. This work is shared under the Creative Commons 
    Attribution-NonCommercial-ShareAlike (CC BY-NC-SA) license, allowing 
    internal use and adaptation while respecting the original creators' rights.
    % \small Key Permissions under CC BY-NC-SA 4.0:
    % \begin{itemize}
    %   \item credit the creator.
    %   \item no commercial use.
    %   \item new creations must carry the same license.
    %   \item copy and redistribute the material in any medium or format
    %   \item remix, transform, and build upon the material
    % \end{itemize}
    }\\[0.5cm]

    % Some content may be derived from sources under the CC0 1.0 Universal 
    % license, which allows for free use, modification, and distribution.}\\[0.5cm]
    {\small\textbf{Quality Assurance:}}\\[0.2cm]
    {\footnotesize
    We have carefully reviewed this material for accuracy and clarity. However, 
    as with any educational resource, we encourage critical engagement and 
    verification of concepts. If you notice any issues or have suggestions for 
    improvement, please bring them to our attention.}
  \end{minipage}
  \includegraphics[width=0.45\textwidth]{images/CA_LastVersion.png}\\
  {\small © \the\year\ Hans edu LLC. All rights reserved.}\\[0.5cm]
  {\small Written by \BookAuthor}\\[1cm]
\end{center}

\newpage


\begin{enumerate}
  \item \textbf{Arc Length Calculation} (10 points)\\
  \insertimage{0.30}{images/2025_06_15_44388a8c3a9d04fd81aag-01}{Circle with center O and arcs}
  The circle above has center $O$, the length of arc $\overline{ADC}$ is $5\pi$, and $x=100$. What is the length of arc $\overline{ABC}$?
  \begin{enumerate}[label=(\Alph*)]
    \item $9\pi$
    \item $13\pi$
    \item $18\pi$
    \item $\frac{13}{2}\pi$
  \end{enumerate}
  \begin{subanswer}
    % your answer here
  \end{subanswer}

  \item \textbf{Circle Radius from Equation} (10 points)\\
  The graph of $x^{2}+x+y^{2}+y=\frac{199}{2}$ in the $xy$-plane is a circle. What is the length of the circle's radius?
  \begin{subanswer}
    % your answer here
  \end{subanswer}

  \item \textbf{Circle Center Coordinates} (10 points)\\
  The equation above defines a circle in the $xy$-plane. What are the coordinates of the center of the circle?
  \begin{enumerate}[label=(\Alph*)]
    \item $(-20,-16)$
    \item $(-10,-8)$
    \item $(10,8)$
    \item $(20,16)$
  \end{enumerate}
  \begin{subanswer}
    % your answer here
  \end{subanswer}

  \newpage

  \item \textbf{Arc Length with Angle} (10 points)\\
  \insertimage{0.40}{images/2025_06_15_44388a8c3a9d04fd81aag-04}{reference attached}
  Point $O$ is the center of the circle above, and the measure of $\angle OAB$ is $30^{\circ}$. If the length of $\overline{OC}$ is 18, what is the length of arc $\overline{AB}$?
  \begin{enumerate}[label=(\Alph*)]
    \item $9\pi$
    \item $12\pi$
    \item $15\pi$
    \item $18\pi$
  \end{enumerate}
  \begin{subanswer}
    % your answer here
  \end{subanswer}

  \item \textbf{Circle Radius from Diameter} (10 points)\\
  A circle in the $xy$-plane has a diameter with endpoints $(2,4)$ and $(2,14)$. An equation of this circle is $(x-2)^{2}+(y-9)^{2}=r^{2}$, where $r$ is a positive constant. What is the value of $r$?
  \begin{subanswer}
    % your answer here
  \end{subanswer}

  \newpage

  \item \textbf{Angle Calculation in Triangle} (10 points)\\
  \insertimage{0.40}{images/2025_06_15_d0312806c0eedc278ae1g-01}{reference attached}
  In the figure, $AC=CD$. The measure of angle $EBC$ is $45^{\circ}$, and the measure of angle $ACD$ is $104^{\circ}$. What is the value of $x$?
  \begin{subanswer}
    % your answer here
  \end{subanswer}

  \item \textbf{Complex Angle Calculation} (10 points)\\
  \insertimage{0.40}{images/2025_06_15_d0312806c0eedc278ae1g-02}{reference attached}
  In the figure shown, points $Q$, $R$, $S$, and $T$ lie on line segment $PV$, and line segment $RU$ intersects line segment $SX$ at point $W$. The measure of $\angle SQX$ is $48^{\circ}$, the measure of $\angle SXQ$ is $86^{\circ}$, the measure of $\angle SWU$ is $85^{\circ}$, and the measure of $\angle VTU$ is $162^{\circ}$. What is the measure, in degrees, of $\angle TUR$?
  \begin{subanswer}
    % your answer here
  \end{subanswer}

  \newpage

  \item \textbf{Intersecting Lines Angle} (10 points)\\
  Intersecting lines $r$, $s$, and $t$ are shown below.
  \insertimage{0.35}{images/2025_06_15_d0312806c0eedc278ae1g-03}{reference attached}
  What is the value of $x$?
  \begin{subanswer}
    % your answer here
  \end{subanswer}

  \item \textbf{Parallel Lines Angle} (10 points)\\
  \insertimage{0.35}{images/2025_06_15_d0312806c0eedc278ae1g-04}{reference attached}
  In the figure, line $m$ is parallel to line $n$. What is the value of $w$?
  \begin{enumerate}[label=(\Alph*)]
    \item 17
    \item 30
    \item 70
    \item 170
  \end{enumerate}
  \begin{subanswer}
    % your answer here
  \end{subanswer}

  \newpage

  \item \textbf{Isosceles Triangle Angle} (10 points)\\
  \insertimage{0.20}{images/cus07.png}{reference attached}
  In the given triangle, $AB=AC$ and $\angle ABC$ has a measure of $67^{\circ}$. What is the value of $x$?
  \begin{enumerate}[label=(\Alph*)]
    \item 36
    \item 46
    \item 58
    \item 70
  \end{enumerate}
  \begin{subanswer}
    % your answer here
  \end{subanswer}

  \item \textbf{Intersecting Segments Angle} (10 points)\\
  \insertimage{0.40}{images/2025_06_15_d0312806c0eedc278ae1g-06}{reference attached}
  
  In the figure above, $\overline{MQ}$ and $\overline{NR}$ intersect at point $P$, $NP=QP$, and $MP=PR$. What is the measure, in degrees, of $\angle QMR$? (Disregard the degree symbol when gridding your answer.)
  \begin{subanswer}
    % your answer here
  \end{subanswer}


  \newpage

  \item \textbf{Similar Triangles Angle} (10 points)\\
  \insertimage{0.35}{images/cus08.png}{reference attached}
  Right triangles $PQR$ and $STU$ are similar, where $P$ corresponds to $S$. If the measure of angle $Q$ is $18^{\circ}$, what is the measure of angle $S$?
  \begin{enumerate}[label=(\Alph*)]
    \item $18^{\circ}$
    \item $72^{\circ}$
    \item $82^{\circ}$
    \item $162^{\circ}$
  \end{enumerate}
  \begin{subanswer}
    % your answer here
  \end{subanswer}

  \item \textbf{Parallel Lines with Transversal} (10 points)\\
  \insertimage{0.35}{images/2025_06_15_d0312806c0eedc278ae1g-08}{reference attached}
  In the figure, line $q$ is parallel to line $r$, and both lines are intersected by line $s$. If $y=2x+8$, what is the value of $x$?
  \begin{subanswer}
    % your answer here
  \end{subanswer}

  \newpage

  \item \textbf{Parallel Lines Angle Relationship} (10 points)\\
  \insertimage{0.35}{images/2025_06_15_d0312806c0eedc278ae1g-09}{reference attached}
  In the figure above, lines $\ell$ and $m$ are parallel, $y=20$, and $z=60$. What is the value of $x$?
  \begin{enumerate}[label=(\Alph*)]
    \item 120
    \item 100
    \item 90
    \item 80
  \end{enumerate}
  \begin{subanswer}
    % your answer here
  \end{subanswer}

  \item \textbf{Parallel Lines Proportionality} (10 points)\\
  \insertimage{0.30}{images/2025_06_15_d0312806c0eedc278ae1g-10}{reference attached}
  In the figure above, $\overline{AF}$, $\overline{BE}$, and $\overline{CD}$ are parallel. Points $B$ and $E$ lie on $\overline{AC}$ and $\overline{FD}$, respectively. If $AB=9$, $BC=18.5$, and $FE=8.5$, what is the length of $\overline{ED}$, to the nearest tenth?
  \begin{enumerate}[label=(\Alph*)]
    \item 16.8
    \item 17.5
    \item 18.4
    \item 19.6
  \end{enumerate}
  \begin{subanswer}
    % your answer here
  \end{subanswer}

  \newpage

  \item \textbf{Similar Triangles Sine Value} (10 points)\\
  Triangle $FGH$ is similar to triangle $JKL$, where angle $F$ corresponds to angle $J$ and angles $G$ and $K$ are right angles. If $\sin(F)=\frac{308}{317}$, what is the value of $\sin(J)$?
  \begin{enumerate}[label=(\Alph*)]
    \item $\frac{75}{317}$
    \item $\frac{308}{317}$
    \item $\frac{317}{308}$
    \item $\frac{317}{75}$
  \end{enumerate}
  \begin{subanswer}
    % your answer here
  \end{subanswer}

  \item \textbf{Right Triangle Trigonometric Relationship} (10 points)\\
  In right triangle $RST$, the sum of the measures of angle $R$ and angle $S$ is 90 degrees. The value of $\sin(R)$ is $\frac{\sqrt{15}}{4}$. What is the value of $\cos(S)$?
  \begin{enumerate}[label=(\Alph*)]
    \item $\frac{\sqrt{15}}{15}$
    \item $\frac{\sqrt{15}}{4}$
    \item $\frac{4\sqrt{15}}{15}$
    \item $\sqrt{15}$
  \end{enumerate}
  \begin{subanswer}
    % your answer here
  \end{subanswer}

  \newpage

  \item \textbf{Tangent Value Calculation} (10 points)\\
  \insertimage{0.30}{images/cus09.png}{reference attached}
  In the figure above, what is the value of $\tan(A)$?
  \begin{enumerate}[label=(\Alph*)]
    \item $\frac{20}{29}$
    \item $\frac{21}{29}$
    \item $\frac{20}{21}$
    \item $\frac{21}{20}$
  \end{enumerate}
  \begin{subanswer}
    % your answer here
  \end{subanswer}

  \item \textbf{Right Triangle Side Length} (10 points)\\
  \insertimage{0.15}{images/cus10.png}{reference attached}
  Triangle $ABC$ above is a right triangle, and $\sin(B)=\frac{5}{13}$. What is the length of side $\overline{BC}$?
  \begin{subanswer}
    % your answer here
  \end{subanswer}

  \newpage

  \item \textbf{Pythagorean Theorem Application} (10 points)\\
  \insertimage{0.30}{images/cus11.png}{reference attached}
  Which equation shows the relationship between the side lengths of the given triangle?
  \begin{enumerate}[label=(\Alph*)]
    \item $4b=19$
    \item $4+b=19$
    \item $4^{2}+b^{2}=19^{2}$
    \item $4^{2}-b^{2}=19^{2}$
  \end{enumerate}
  \begin{subanswer}
    % your answer here
  \end{subanswer}
\end{enumerate}